% Long-only authored mathematics for the Erdős #68 reasoning record.
% Everything here was in the short note before the September 2026 revision and
% is kept because the complete record owns the subsidiary arguments, the
% superseded deductions, and the failed routes with their exact boundaries.

\section{Extended record}\label{long68:sec:extended-record}

The short note above carries the advertised results with complete proofs.
This part keeps the surrounding mathematics: the digit kernel the carry and
orbit criteria run on, the exact interfaces of the collision-core reduction,
the finite certificates, the superseded deductions, and the routes that were
tried and closed.

\subsection{The canonical factorial digit kernel}\label{long68:sec:ext-digits}

Write $\theta_0=\{x\}$ and, for $m\ge1$,
\[
 a_m=\lfloor m\,\theta_{m-1}\rfloor,\qquad \theta_m=m\,\theta_{m-1}-a_m ,
\]
so that $\theta_m=\{m!\,x\}$ for $m\ge1$.  The kernel checks the floor
formula, the digit bounds $0\le a_m<m$, the recurrence
$\theta_{m+1}=(m+1)\theta_m-a_{m+1}$, the finite telescoping expansion
\[
 x=\lfloor x\rfloor+\sum_{m=2}^{N}\frac{a_m}{m!}+\frac{\theta_N}{N!},
\]
and the rule that a zero remainder at one index forces every later digit to
vanish.  The rational direction is also checked: if $q>0$ and $q\le n$, then
$\operatorname{facFloor}(a/q,n)=((n!/q):\Z)\,a$ and the canonical digit at
radix $n+1$ vanishes, so every rational input has an eventually zero
expansion.  The converse holds for every real input, and gives
\[
 x\in\Q\quad\Longleftrightarrow\quad a_m(x)=0\ \hbox{for all large }m .
\]
If all digits after index $N\ge1$ vanish, the recurrence gives
$\theta_{N+k}\ge(k+1)\theta_N$; since every remainder is below one,
$\theta_N=0$ and $x=\lfloor N!x\rfloor/N!$.

A zero canonical digit is a different event from a zero-branch hit.  The
returned data contain canonical zero digits at $m=5$ and $m=23$, while the
zero-branch list is empty through $m=100000$.

A second exact reformulation runs through a defect automaton.  For a rational
centre recurrence $F_m=mF_{m-1}+1+\varepsilon_m-C_m$, the kernel checks that
the integer ceiling defect code equals
$\lfloor m\delta_{m-1}-\varepsilon_m\rfloor$ and that
$\delta_m=m\delta_{m-1}-\varepsilon_m-q_m$, with the specialisation
$\varepsilon_m=1/(m!-1)$ written out.  What is checked is the algebra of the
automaton; proving that the finite-sum residual centre satisfies the premise
is a separate step and is not done.

\subsection{Criteria from the literature, in full}\label{long68:sec:ext-literature}

Koepf and Schmersau prove that eventual equality between the floors of $n$
times a partial sum and $n$ times its limit forces irrationality
\cite[Thm.~1.1, p.~117]{koepf-schmersau}; their rational-term version obtains
that equality from prefix integrality at a scale $p_n$ together with the
strict tail bound $a-s_n<1/(np_n)$
\cite[Thms.~2.2--2.3, pp.~119--120]{koepf-schmersau}.  For the natural
termwise clearing choice $p_n=\lcm\{k!-1:2\le k\le n\}$ the last two
denominators already obstruct it:
\[
 p_n\ \ge\ \frac{(n!-1)\bigl((n-1)!-1\bigr)}{n-1},
\]
because $\gcd(n!-1,(n-1)!-1)=\gcd((n-1)!-1,n-1)\le n-1$.  For $n\ge4$ the
first omitted summand $1/((n+1)!-1)$ then already exceeds $1/(np_n)$, so this
$p_n$ cannot satisfy their tail hypothesis.  Cancellation in the reduced
prefix denominator could give a smaller scale, and proving enough
cancellation is another form of the present denominator problem.

Duverney's Theorem~3.1 assumes two-sided quadratic denominator growth
$cu_n^2\le u_{n+1}\le c'u_n^2$, while $u_{n+1}/u_n^2\to0$ for $u_n=n!-1$
\cite[pp.~275, 285--286]{duverney}; his printed Corollary~3.2 displays the
signed series $\sum_n(u_{n+1}/u_n^2-1)$, without absolute values.  Its
summands tend to $-1$ here, so it does not converge to a finite value
\cite[Cor.~3.2, p.~287]{duverney}.  The ratio proof divides numerator
and denominator by $(n!)^2$; the resulting terms tend to zero while
$(1-1/n!)^2$ tends to one.

The theorem of Barreto, Kang, Kim, Kova\v{c} and Zhang proves irrationality
of $\sum_n1/a_n$ when $a_n^{1/2^n}\to\infty$, whereas $(n!-1)^{1/2^n}\to1$
\cite[Thms.~2--3, pp.~2--4]{barreto-et-al}; the limit is kernel-checked
through the bound $0\le\log(n!-1)\le n^2$.  Their proof identifies a useful
exact criterion: Mahler's elementary rationality floor is contradicted by
prefix-clearing integers $D_N$ whose cleared positive tails satisfy
$\liminf_ND_Nr_N=0$ \cite[Lem.~8 and Prop.~12, pp.~6, 9--12]{barreto-et-al}.
The ordinary product of the factorial-gap denominators is far too large for
that estimate, so a transfer needs a low-height clearing subsequence, or
enough exact cancellation in their least common multiple.

Dividing the strict-successor recurrence $Z_m=mZ_{m-1}+1-b_m$ by $m!$ and
telescoping gives the exact finite identity
\[
 \frac{Z_M}{M!}=\frac{Z_2}{2!}+\sum_{m=3}^{M}\frac{1-b_m}{m!},
\]
so the carry defects $1-b_m$ are genuine factorial-series coefficients.
Han\v{c}l and Tijdeman give exact rationality classifications for polynomial
coefficients and finite-difference criteria for broader ordinary factorial
series \cite[Thm.~3.1 and Cor.~3.1, pp.~390--391]{hancl-tijdeman}; their
denominator is the cumulative linear product $\prod_{n\le N}(an+b)$, and the
individual number $N!-1$ does not occur.  Applied to the display above, the Cantor--Oppenheim criterion
still needs $1-b_m\ne0$ infinitely often, which is the missing cofinal
assertion.

\subsection{A superseded deduction}\label{long68:sec:ext-superseded}

Before the elementary segment argument of the short note, the record derived
a weaker exponent from an external multiplicity theorem.  Put
$Q_N=\prod_{2\le n\le N}(n!-1)$.  For an odd prime $r$, let $m_r$ count the
indices $2\le n\le N$ with $r\mid n!-1$; such an index satisfies $n<r$.  The
factorial-congruence multiplicity estimate of Garaev, Luca and Shparlinski
\cite[arXiv v1, Thm.~12, p.~16]{garaev-luca-shparlinski}, applied on
$1\le n\le\min(N,r-1)$, gives $m_r\ll N^{2/3}$; the prime $2$ divides none of
these factors.  Writing $E_r=\max_{2\le n\le N}v_r(n!-1)$,
\[
 \log Q_N=\sum_r\sum_{n=2}^{N}v_r(n!-1)\log r
 \le\bigl(\max_r m_r\bigr)\sum_rE_r\log r
 \ll N^{2/3}\log \Lc_N ,
\]
and Stirling summation gives $\log Q_N\asymp N^2\log N$, whence
$\log \Lc_N\gg N^{4/3}\log N$.  Theorem~\ref{long68:res:lcm-growth} supersedes this:
its exponent is $3/2$, its constant is explicit, and it uses no external
theorem.  The deduction is retained because it is the shape the returned
analyses describe, and because it is the only place a multiplicity bound
enters the record.

The primitive lcm divisibility survives cofactor removal: factorial
valuations do not remove the channel obstruction once every common cofactor
divisor has been removed.  A corank-one cofactor and determinant construction
for such primitive kernels has not been formalised, and the divisibility
theorem does not establish it.

\subsection{Plateaux, first exit, and the finite denominator ladder}
\label{long68:sec:ext-plateau}

A second argument works on the rational grid in place of the channels.  Let
$H$ be a partial sum and $q$ a candidate denominator.  Writing $qH=k+r$ and
$q(\Ser-H)=u$, the next $q^{-1}$ grid point $(k+1)/q$ lies below $\Ser$
exactly when $1\le r+u$.  The factorial plateau theorem states that if
$H<G\le\Ser$, if $n!G$ is integral, and if $n!(\Ser-H)<1$, then the strict
successor of $n!H$ and the canonical floor of $n!\Ser$ are the same grid
integer.  Two rigidity statements follow: consecutive plateau floors, scaled
by the next radix, force the canonical digit to vanish; and any first-exit
offset $\delta\in[0,2)$ with carry $b=-\lfloor\delta\rfloor$ has
$b\in\{0,-1\}$, so the exit has exactly two alternatives.

The first-crossing argument continues to a denominator bound.  For a rational
grid level $G$ with first crossing $\tau$ by the literal partial sums, write
$G-H_{\tau-1}=a/v$ with $a,v>0$.  Then $v\ge\tau!-1$, and on the $-1$ exit
branch $v\ge\tau!(\tau!-1)$.  No coprimality hypothesis on $a$ and $v$ is
needed.

There is also a direct obstruction at prime indices.  With $\Delta_m$ as in
the short note, the kernel checks for $m\ge3$ the criterion
\[
 m\mid Z_m
 \quad\Longleftrightarrow\quad
 1+\frac1{m!-1}<m\Delta_m\le2+\frac1{m!-1},
\]
and if $\Ser=a/q$ with $q>0$, then $p\mid Z_p$ for every prime $p>q$, so one
exact missed prime $p$ gives $q\ge p$.  Exact kernel reduction gives
$11\nmid Z_{11}$, hence $q\ge11$; exact rational normalisation gives
$60\nmid Z_{60}$, $64\nmid Z_{64}$ and $67\nmid Z_{67}$, and the prime index
$67$ gives the checked bound
\[
 \Ser=\frac aq,\ q>0\quad\Longrightarrow\quad q\ge67 .
\]
The exact-interval census of the short note replaces $67$ by $300000$.

There is a finite peeling identity.  For $x\ne0,1$ and $K\ge0$,
\[
 \frac1{x-1}=\sum_{j=1}^{K}\frac1{x^j}+\frac1{x^K(x-1)} .
\]
When $x=k!$ and a chosen factorial scale is divisible by $(k!)^K$, the scaled
finite sum is integral and only the last term retains the factor $k!-1$ in
its denominator.  The identity isolates one residual fraction before exact
bounding, and it supplies no cofinal family of nonzero residuals.

One proposed strengthening is false and is recorded as false: the
divisibility $(m!-1)\mid\den(V_m)$ fails at the reported strict events
$m=52$ and $m=591$.  Only the Archimedean first-crossing lower bound
survives.

There is a second, more arithmetic mechanism at doubled prime indices, stated
as \eqref{long68:eq:doubled-prime} in the short note.  The formal theorem is not
restricted to hand-reduced indices; it is the general unique-slot prime-power
criterion specialised to the literal prefixes.

\subsection{The collision-core interfaces}\label{long68:sec:ext-collision}

Beyond the definitions used in the short note, the record keeps the exact
interfaces of the collision-core reduction.

The collision core has an exact incremental law.  For the positive
factorial-gap denominators, adjoining $d_a$ to a finite family $S$ gives
\[
 C(S\cup\{a\})=\lcm\!\Bigl(C(S),\ \gcd\bigl(d_a,\lcm_{j\in S}d_j\bigr)\Bigr),
\]
since finite-family gcd and lcm distributivity collapses the lcm of all
pairwise gcds against $d_a$ to a single gcd.  The same formula holds after
adjoining the distinguished base, so each step needs only the old denominator
lcm and the old core, with no pairwise rescan.

There is an exact product and lcm bound.  If $\widetilde C(S)$ is the core
after cancelling a positive distinguished base, $L(S)=\lcm_{j\in S}d_j$ and
$P(S)=\prod_{j\in S}d_j$, then $\widetilde C(S)L(S)\mid P(S)$, hence
$\widetilde C(S)\le P(S)/L(S)$.  For the factorial block this specialises to
\[
 \widetilde C_p\le\frac{\prod_{n\in I_p}(n!-1)}{\lcm_{n\in I_p}(n!-1)},
\]
an exact bridge from lower estimates for the factorial-gap lcm to upper
estimates for the normalised core.  It does not close the local scale bound.

The base cancellation is exact prime by prime.  With $F_p=(p-1)!$ and $C_p$
the unnormalised core,
\[
 \widetilde C_p=\frac{\lcm(F_p,C_p)}{F_p}=\frac{C_p}{\gcd(F_p,C_p)},\qquad
 v_r(\widetilde C_p)=v_r(C_p)-\min\{v_r(F_p),v_r(C_p)\} .
\]
So $r^e\mid\widetilde C_p$ exactly when the pairwise core carries
$r^{e+v_r(F_p)}$, which in the block forces two distinct gaps to be divisible
by that higher power.  The sharp surviving valuation cap
$v_r(\widetilde C_p)+v_r(F_p)<r$ then gives $p-1<r(r-1)<r^2$ for every
support prime, and $\widetilde C_p$ is coprime to $k!$ whenever
$k(k-1)\le p-1$.  This removes every factorial channel below the moving
square-root cutoff, and it does not bound the aggregate product of the
remaining large prime powers.

For collision estimates that already provide an upper-half hit, no exponent
is lost to normalisation: if $r$ divides a displayed gap at some $n\ge p$,
then $r\nmid F_p$ and $r^e\mid\widetilde C_p\iff r^e\mid C_p$ for every
$e>0$.  This has an exact incidence-count form,
\[
 r^e\mid\widetilde C_p
 \quad\Longleftrightarrow\quad
 1<\#\{i\in I_p:r^e\mid i!-1\},
\]
so a source estimate giving at most one $r^e$-hit deletes that exponent and
yields $v_r(\widetilde C_p)<e$.  The local aggregation is exact:
\[
 v_r(\widetilde C_p)=\#\bigl\{e\in[1,r-1]:1<\#\{i\in I_p:r^e\mid i!-1\}\bigr\},
\]
and for an endpoint prime with $2p-1<r$ the range truncates to
$e\in[1,2p-4]$.  The same equivalence holds under the exact condition
$r\nmid F_p$ in place of the endpoint inequality, which covers the entire
moving prime range at and above the block parameter; at $e=2$ it gives a
conditional squarefree conclusion $v_r(\widetilde C_p)\le1$.

The local load has a distance-sensitive witness.  If $r$ is prime, $e>0$ and
$r^e\mid\widetilde C_p$, there are $i<j$ in $I_p$ with
$r^{e+v_r(F_p)}\mid i!-1$, $r^{e+v_r(F_p)}\mid j!-1$ and
$r^{e+v_r(F_p)}\le j^{\,j-i}$; consequently $(2p-1)^d<r^{e+v_r(F_p)}$ forces
$d<j-i$.  The spacing hypothesis is discharged internally: any two $r^e$-hits
$i<j$ satisfy $e<j-i$, because $r\mid j!-1$ already forces $j<r$ and the
gap-power inequality converts that size relation into strict separation.  The
pairwise form is stronger than the selected-witness form: for arbitrary $r,e$
and any displayed hits $i<j$,
\[
 r^e\mid i!-1,\quad r^e\mid j!-1
 \quad\Longrightarrow\quad r^e\le j^{\,j-i},
\]
so every prime-power hit layer is an $e$-separated subset of the block and
\[
 (e+1)\,\#\{i\in I_p:r^e\mid i!-1\}\ \le\ 2p+e-2 .
\]
The unweighted packing step is therefore complete.  Globally,
\[
 r\mid\widetilde C_p
 \quad\Longrightarrow\quad
 v_r(\widetilde C_p)+v_r(F_p)<2p-3
\]
for every prime $r$ and $p\ge2$, so even primes already present in the
normalisation base pay for their base valuation inside the same
block-diameter budget.  What remains is to bound the layer counts uniformly
as $r$ and $p$ move, multiply the surviving contributions, and still close
the complementary-residue coordinate.

The squarefreeness premise is not proved.  An exhaustive modular scan through
$r\le2{,}000{,}000$ and $n\le240$ found four individual square hits and no
prime with two such hits.  All $498{,}501$ pairs $2\le a<b\le1000$ have
squarefree $\gcd(a!-1,b!-1)$, and the aggregate squarefree-collision scan
through $p=499$ stays below $0.374$ of the upper-descending-factorial
logarithmic scale.  These are finite exact computations.  They carry neither theorem authority
nor an asymptotic incidence bound.

A finite variant of the cofinal private supply compares the product of a
chosen set of primes, each at least $5$, with $\prod_{2\le k\le B}(k!-1)$; if
the prime product is larger, at least one chosen prime has no hit through
$B$, while Wilson still bounds its least hit by $q-2$.  Wilson reflection
limits what a linear-size divisor can be assumed to be: if $n$ is odd,
$n<q$, and $q\mid n!-1$, then $q\mid(q-n-1)!-1$, and when both indices lie in
one block with the reflected hit earlier, equivalently $q<2n+1$, the repeated
hit survives predecessor-factorial normalisation and its full-block incidence
count exceeds one.

Stewart states that for every $\varepsilon>0$ there are infinitely many odd
$n$ whose least prime factor $q$ of $n!-1$ satisfies
\[
 n<q<\left(\frac{\sqrt{145}-1}{8}+\varepsilon\right)n ,
\]
the printed text transferring estimate~(9) from $n!+1$ to $n!-1$
\cite[p.~464]{stewart2004}.  The source controls $q$ relative to the later
index $n$ and supplies no control relative to the private first-hit index,
so it is collision-core input.  The missing private-anchor and global product
estimates are elsewhere.

One selected prime $q$ already furnishes the coprime factor pair $(1,q)$: its
projection moduli are $R_p$ and $R_p/q$, whose least common multiple is
$R_p$, so no second selected prime is needed and there is no hidden
equality-or-disagreement branch in that specialisation.  The underlying
projection rigidity is general: if the endpoint numerator $Z$ is congruent to
a weighted numerator $T$ modulo $R$, then every divisor $Q$ of $R$ with
$Z\le B<Q$ satisfies $T\bmod Q=Z$, so two divisors above $B$ with unequal
projected residues exclude that bounded endpoint, and unequal projections
force $\min(Q_1,Q_2)\le T$.

\subsection{Finite channel and carry certificates}\label{long68:sec:ext-certificates}

The finite-support vector $\lambda=2e_3-e_4$ has, by kernel check,
$\chn2\lambda=0$, factorial moment $-12$, $\chn3\lambda=-2$,
$\chn4\lambda=11$, and $\chn d\lambda=-12$ for every $d\ge5$.  Under the
exact rational tail enclosure $1/119<\Theta_4<1/50$ its residual lies
strictly between $-93/575$ and $-309/13685$, so it is nonzero and subunit.

Exact integer regeneration verifies the canonical primitive kernels for every
$2\le D\le12$: channels $2$ through $D$ vanish, the factorial moment is
$\Lc_D$, and the coefficient content is one.  At $D=9$ the moment is
$\Lc_9=31540008254514077395$ and, after the stated prime-unit shift,
$1353/100000<R_9<1354/100000$.  At $D=3$ the vector $c=(-40,55,-10,1)$ on the
support $(3,4,5,6)$ annihilates channels $2$ and $3$, has moment $600$, and
satisfies
\[
 0.09925341997208298<\Lc_3(c)<0.09925341997208300 ,
\]
which excludes denominators dividing $600$.

There are two different computations at the same endpoint.  A returned
interval computation reports the stronger geometric statement that no
zero-branch event occurs at any $m\le100000$; its cited executable and source
digest were not supplied, so that classification remains external finite
evidence.  The strict-successor carry computation used in the short note is
local and independently regenerated, with its own source, backend, payload
and receipt digests.

\subsection{Limits of fixed-coordinate arguments}\label{long68:sec:ext-nogo}

\begin{itemize}
\item Residue vectors, their recurrences, and window widths admit synthetic
  all-hit blocks, so they cannot prove irrationality on their own.
\item Known pointwise prime congruences, prime-dilation congruences, parity,
  and the exact prime coefficient formula admit a synthetic rational
  countermodel.
\item Wilson quotients, harmonic sums, $p$-adic gamma identities, and
  factorial residues do not control the required Archimedean floor without an
  additional coupling theorem.  Every prime-window test factors into a sharp
  Archimedean strict-ceiling condition and a modular divisibility condition,
  and the missing ingredient is the coupling between them.
\item Fixed-denominator scalar canonical-product localisers and
  rank-saturated consecutive-jet Hermite--Pad\'e systems pay the full
  factorial-gap denominator.  For $E(z)=\prod_{n\ge2}(1-z/n!)$ the genus-zero
  product satisfies $-E'(1)/E(1)=\Ser$, and the natural scalar linear form
  carries the coefficient $Q_N=\prod_{2\le n\le N}(n!-1)$, for which $Q_N$
  times the tail diverges.  Exact first-order interpolation, scalar residue
  weighting, the natural Wronskian, and rank-saturated consecutive jets all
  reassemble the same prohibitive denominator.
\item Zero-moment variations cannot create an additional fractional
  cancellation coordinate, and factorial valuations cannot absorb the channel
  lcm obstruction.
\item A fixed pair of low-index private owners cannot make the projection
  argument cofinal, by the absorption statement of the short note.
\end{itemize}

\subsection{The generic shape behind the lower-interval form}
\label{long68:sec:ext-generic}

The cofinal statement in the short note has a generic ancestor.  For a real
$x$ and any positive sequence $c_m$ with $c_m<1/m$ eventually,
\[
 x\notin\Q
 \quad\Longleftrightarrow\quad
 \{(m-1)!\,x\}\ge c_m\ \hbox{for arbitrarily large }m .
\]
For rational $x$ the fractional parts vanish eventually; if all large
fractional parts lie below $c_m$, multiplying by $m$ stays below one, so the
next canonical digit is zero, and eventual zero digits force rationality.
The specific target $c_m=E_m/m$ is useful because it translates into finite
predecessor-gap conditions.  The smallness of the target alone establishes no
metric, measure or equidistribution theorem for this orbit, and none is
claimed.
